\section{Introduction}

Housing markets play a central role in household wealth, local economies, and monetary policy. According to Matteo Iacoviello, housing wealth represents ``one half of household net worth'' \cite{iacoviello2011}, which makes it a determinant of the broader economy: losses in housing wealth can dramatically decrease consumption, with consequences for monetary policy and GDP growth. Because housing amplifies economic expansions and recessions, predicting housing prices in real time has real practical value. However, official indicators such as the Case-Shiller Index and the Federal Housing Finance Agency price index are consistently published with latency, which limits the ability of policymakers, investors, and economists to assess current market conditions.

This work addresses that gap empirically. Nowcasting refers to a framework that attempts to ``estimate complex case counts for a given reporting date, using time series of case reports that are known to be incomplete due to reporting delays'' \cite{mcgough2020}; in our setting it means using real-time indicators to predict the current month's housing price before the official Zillow figure is released. We do not propose a new method. Rather, our goal is to \emph{benchmark} four standard predictive approaches --- a pooled OLS regression with metropolitan fixed effects, a univariate ARIMA model, a gradient-boosted tree model (XGBoost), and a two-layer neural network (MLP) --- on three structurally different US metropolitan areas: San Francisco, California; Austin, Texas; and Cleveland, Ohio.

The questions we ask are descriptive and comparative. \textbf{(i)~Do specific features emerge as the major drivers of monthly housing-price returns, or does the answer depend on the metropolitan area?} \textbf{(ii)~Is model selection itself geographically contingent in the United States --- that is, does the best-performing model differ across structurally different metros, or does one approach dominate everywhere?} The models are trained on pre-2023 data and evaluated on a January~2023--March~2026 test set that covers the post-2022 mortgage-rate-shock regime, providing a demanding out-of-sample comparison.
